% 3. Bernstein-Zelevinsky（及其他）分类（底本 p.23--31）
\section{Bernstein--Zelevinsky（及其他）分类}

\subsection{$\GL_n$ 容许表示的分类}

\subsubsection{段 $\Delta(\pi,m)$}

设 $(\pi,V)\in\Repadm_G$，对任意 $s\in\mathbb{C}$，记 $\pi(s)$ 为 $\pi$ 用特征标
$|\det|^{s}$ 所作的扭转，即 $\pi(s)(g)=|\det(g)|^{s}\pi(g)$。
注意若 $\pi$ 不可约超尖点，则 $\pi(s)$ 也是。

我们在 $\GL_n(K)$ 的不可约超尖点表示的等构类集合上给一个偏序：
$\pi\le\pi'$ 当且仅当存在整数 $s\ge 0$ 使 $\pi'=\pi(s)$。记这个偏序集为 $\Asc_n$。
$\Asc_n$ 中的有限区间称为\emph{Zelevinsky 段}。更精确地说：

\begin{definition}
从 $\pi$ 出发、长度为 $m$、次数为 $mn$ 的\emph{Zelevinsky 段}\en{9}是
$\GL_n(K)$ 的 $m$ 个不可约超尖点表示构成的有序有限族：
\[
\Delta(\pi,m)=[\pi,\pi(1),\dots,\pi(m-1)].
\]
\end{definition}

\begin{definition}
设 $\Delta_1=(\pi_1,m_1)$ 与 $\Delta_2=(\pi_2,m_2)$ 分别是 $\Asc_{n_1}$ 与 $\Asc_{n_2}$
中的两个 Zelevinsky 段。
\begin{enumerate}[label=(\arabic*),leftmargin=2.2em]
\item 若 $\Delta_1\not\subseteq\Delta_2$、$\Delta_2\not\subseteq\Delta_1$，且
      $\Delta_1\cup\Delta_2$ 是 Zelevinsky 段，则称 $\Delta_1$ 与 $\Delta_2$ \emph{相链}
      （linked）。特别地，这只有在 $n_1=n_2$ 时才可能发生。
\item 若 (a) $\Delta_1$ 与 $\Delta_2$ 相链，且 (b) $\pi_1<\pi_2$，则称
      $\Delta_1$ \emph{先于} $\Delta_2$（precedes）。
\end{enumerate}
\end{definition}

\begin{remark}
特别地，若 $n_1\neq n_2$，则 $\Delta_1$ 与 $\Delta_2$ 不相链。
\end{remark}

\begin{definition}
设 $\Delta=\Delta(\pi,m)$ 是（Zelevinsky）段，长度为 $m$、次数为 $mn$，
定义 $\GL_{mn}$ 的容许表示
\[
\pi(\Delta):=\pi_1\times\cdots\times\pi_m,
\]
其中 $\pi_1\times\cdots\times\pi_m=I^{\GL_{mn}}_{P_{(n,n,\dots,n)}}(\otimes\pi_i)$，
如 2.2.4 末的注所述。
\end{definition}

\subsubsection{$(*)$ 条件}

设 $\Delta_1,\dots,\Delta_r$ 是 $r$ 个长度分别为 $m_i$、位于 $\Asc_{n_i}$ 中的
Zelevinsky 段构成的有序族。引入下述 $(*)$ 条件：
\[
(*) \qquad \text{对每个 } i<j,\ \Delta_i\ \text{不先于}\ \Delta_j .
\]
$\Delta_1,\dots,\Delta_r$ 的一个 \emph{$(*)$-置换}是指 $\Delta_i$ 上的一个置换 $\sigma$，
使得 $\Delta_{\sigma(1)},\dots,\Delta_{\sigma(r)}$ 仍满足 $(*)$ 条件。

注意：
\begin{enumerate}[label=(\arabic*),leftmargin=2.2em]
\item 按定义，当 $n_i\neq n_j$ 时该条件为空；
\item $(*)$ 条件允许在 $i<j$ 时 $\Delta_j$ 先于 $\Delta_i$。
\end{enumerate}

\begin{example}
设 $\Delta=[\pi,\pi(1)]$，$\Delta'=[\pi(2),\pi(3)]$，$\Delta''=[\pi(5),\pi(6)]$，
则 $(\Delta,\Delta')$ 不满足 $(*)$，而 $(\Delta',\Delta)$、$(\Delta,\Delta'')$，
以及单个段 $\Delta\cup\Delta'=[\pi,\dots,\pi(3)]$ 都满足 $(*)$。
\end{example}

\subsubsection{分类定理的陈述}

\begin{theorem}[Bernstein--Zelevinsky]
\leavevmode
\begin{enumerate}[label=(\arabic*),leftmargin=2.2em]
\item 对每个长度为 $m$ 的 Zelevinsky 段 $\Delta$，表示 $\pi(\Delta)$ 的合成列长度为 $2^{m-1}$。
\item 设 $\Delta$ 是有限长的 Zelevinsky 段，则 $\pi(\Delta)$ 有唯一的不可约子表示 $Z(\Delta)$
      与唯一的不可约商表示 $Q(\Delta)$。
\item 设 $\Delta_1,\dots,\Delta_r$ 是满足 $(*)$ 条件的 Zelevinsky 段，则表示
\[
Q(\Delta_1)\times\cdots\times Q(\Delta_r)\in\Repadm_{\GL_{\sum n_im_i}}
\]
有唯一的不可约商 $Q(\Delta_1,\dots,\Delta_r)$。此外，在 $(*)$-置换之下，
$Q(\Delta_1,\dots,\Delta_r)$ 与 $\Delta_i$ 的次序无关。
\item $G=\GL_n(K)$ 的任一不可约表示 $(\pi,V)$ 都同构于如下形式：
\[
\pi\cong Q(\Delta_1,\dots,\Delta_r),
\qquad
\Delta_i=(\pi_i,m_i),\quad n=\sum n_im_i,
\]
其中段族 $\Delta_1,\dots,\Delta_r$ 在 $(*)$-共轭之下唯一。
\item 仍如 (2)，设 $\Delta_1,\dots,\Delta_r$ 是满足 $(*)$ 条件的段，则
\[
Q(\Delta_1)\times\cdots\times Q(\Delta_r)
\]
不可约（从而等于 $Q(\Delta_1,\dots,\Delta_r)$）当且仅当任两段 $\Delta_i$ 与 $\Delta_j$ 都不相链。
作为一个特例，我们重新得到 2.2.4 中的注。
\item 对 $\Delta=\Delta(\pi,m)=[\pi,\dots,\pi(m-1)]$，定义
\[
\Delta^{\vee}=[\pi(m-1)^{\vee},\dots,\pi^{\vee}]=\Delta\bigl(\pi^{\vee}(1-m),m\bigr),
\]
则
\[
Q(\Delta_1^{\vee},\dots,\Delta_r^{\vee})=Q(\Delta_1,\dots,\Delta_r)^{\vee}.
\]
\item (3)--(6) 的类似断言对 $Z(\Delta_1)\times\cdots\times Z(\Delta_r)$ 成立：
      它有唯一的不可约子表示 $Z(\Delta_1,\dots,\Delta_r)$。
\end{enumerate}
\end{theorem}

\begin{remark}
任给 Zelevinsky 段的有序族 $\Delta_1,\dots,\Delta_r$，总存在它的某个置换满足 $(*)$ 条件，
于是 $\GL_n(K)$ 的容许表示的等构类与无序族 $\Delta_1,\dots,\Delta_r$（$n=\sum m_in_i$）
一一对应。
\end{remark}

\begin{remark}
可以给各 $\GL_n(K)$ 的容许表示的 Grothendieck 群之直和 $R_K=\bigoplus_{n\ge 0}R_n(K)$
一个乘法结构
\[
([\pi_1],[\pi_2])\longmapsto[\pi_1\times\pi_2],
\]
使 $R_K$ 成为分次交换环。而映射 $Q(\Delta_1,\dots,\Delta_r)\mapsto
Z(\Delta_1,\dots,\Delta_r)$ 定义了 $R_K$ 上的一个对合。
\end{remark}

\subsubsection{超尖点支撑}

\begin{definition}
设 $\pi\cong Q(\Delta_1,\dots,\Delta_r)$ 如 3.1 所述，定义 $\pi$ 的\emph{超尖点支撑}为
下述（无序的）超尖点表示族：
\[
\operatorname{Supp}(\pi)=\bigl\{\pi_i(j)\bigr\}_{i=1,\dots,r;\ j=1,\dots,m_i-1}.
\]
\end{definition}

\begin{remark}[Zelevinsky]
$\pi$ 的超尖点支撑是唯一的一族无序超尖点表示 $\sigma_1,\dots,\sigma_s$，使得
$Q(\Delta_1,\dots,\Delta_r)$ 与 $Z(\Delta_1,\dots,\Delta_r)$ 都是
$\sigma_1\times\cdots\times\sigma_s$ 的子商。
\end{remark}

\subsubsection{Steinberg 表示}

考虑（拟）特征标 $|\cdot|^{\frac{1-n}{2}}\colon K^{\times}\to\mathbb{C}^{\times}$
以及 Zelevinsky 段
\[
\Delta=\Delta\Bigl(|\cdot|^{\frac{1-n}{2}},n\Bigr)
=\Bigl[|\cdot|^{\frac{1-n}{2}},|\cdot|^{\frac{3-n}{2}},\dots,|\cdot|^{\frac{n-1}{2}}\Bigr].
\]
直接计算标准 Borel 子群 $B$ 上的模特征标 $\delta_B$ 在 $T=L_B=\GL_{(1,\dots,1)}(K)$
上的限制。设 $x=\mathrm{Diag}(x_1,\dots,x_n)$ 是对角元为 $x_1,\dots,x_n$ 的对角矩阵，则
\[
\delta(x)=|\det(\ad_U(x))|
=\Bigl|x_1^{n-1}\cdot x_2^{n-3}\cdots x_n^{1-n}\Bigr|.
\]
换言之，对角环面 $T$ 上的表示 $\pi_T$（此时是一个特征标）
$|\cdot|^{\frac{1-n}{2}}\otimes\cdots\otimes|\cdot|^{\frac{n-1}{2}}$ 等于
$\delta^{-\frac12}$。按定义，$\pi(\Delta)$ 有下述简单描述：
\[
\pi(\Delta)=|\cdot|^{\frac{1-n}{2}}\times\cdots\times|\cdot|^{\frac{n-1}{2}}
=\bigl\{\text{光滑函数}\ f\colon B\backslash G\to\mathbb{C}\bigr\}.
\]
因此，唯一的不可约子表示 $Z(\Delta)$ 只能是平凡表示 $\triv$，
它对应于 $B\backslash G$ 上的常值函数。唯一的商 $Q(\Delta)=:\St_n$ 定义为
$\GL_n(K)$ 的（标准）\emph{Steinberg 表示}。它自对偶，即 $\St_n^{\vee}=\St_n$。

\begin{remark}
当 $n=2$、$\Delta=[|\cdot|^{-\frac12},|\cdot|^{\frac12}]$ 时，定理 3.1 的第 (1) 条断言
$\pi(\Delta)$ 长度为 $2$。于是有
\[
0\to\triv\to\pi(\Delta)\to\St_2\to 0 .
\]
\end{remark}

\begin{remark}
另一种做法：考虑 $\Delta_1=[|\cdot|^{\frac12}]$ 与
$\Delta_2=[|\cdot|^{-\frac12}]$，则 $(\Delta_1,\Delta_2)$ 满足 $(*)$ 条件，
且 $\St_2=Z(\Delta_1,\Delta_2)$，$\triv=Q(\Delta_1,\Delta_2)$。作为对比，有短正合列
\[
0\to\triv\to|\cdot|^{-\frac12}\times|\cdot|^{\frac12}\to\St_2\to 0,
\]
\[
0\to\St_2\to|\cdot|^{\frac12}\times|\cdot|^{-\frac12}\to\triv\to 0 .
\]
\end{remark}

\subsubsection{$\GL_2$ 的局部 Langlands 对应}

$n=2$ 的情形格外清楚。在 $\GL_2$ 一侧，$\pi$ 或者超尖点，或者是某个抛物诱导的商
（或者是单个长度为 $2$ 的段，或者是两个长度均为 $1$ 的段）。于是有如下可能：
\begin{enumerate}[label=(\arabic*),leftmargin=2.2em]
\item $\pi$ 是超尖点表示。
\item $\pi=Q(\chi_1,\chi_2)$，其中 $\chi_1/\chi_2\neq|\cdot|^{\pm 1}$，
      于是实际上 $\pi=\chi_1\times\chi_2$——这些称为\emph{主序列表示}。
\item $\pi=Q(\chi|\cdot|^{\frac12},\chi|\cdot|^{-\frac12})$，这是一维表示
      $\triv\otimes\chi$。
\item $\pi=Q(\Delta)$，其中 $\Delta=[\chi|\cdot|^{-\frac12},\chi|\cdot|^{\frac12}]$，
      它等于 $\St_2(\chi):=\St_2\otimes\chi$。
\end{enumerate}
这些称为用 $\chi$ 扭转的 Steinberg 表示（底本作 “Stenberg”，系笔误）。

它们对应的 Weil--Deligne 表示为
\begin{enumerate}[label=(\arabic*),leftmargin=2.2em]
\item 若 $\pi$ 超尖点，则 $\rec_2(\pi)$ 不可约，故
\[
\rec_2(\pi)=(\rho,0),\qquad \rho\ \text{不可约}.
\]
\item 若 $\pi=\chi_1\times\chi_2$ 是主序列、$\chi_1/\chi_2\neq|\cdot|^{\pm 1}$，则
\[
\rec_2(\pi)=(\chi_1\oplus\chi_2,\,0).\en{10}
\]
\item 若 $\pi=\triv\otimes\chi$，则
\[
\rec_2(\pi)=\bigl(\chi\oplus\chi(1),\,0\bigr),
\]
其中 $\chi(1)$ 是 $\chi$ 的「Tate 扭转」。
\item 若 $\pi=\St_2(\chi)$，则
\[
\rec_2(\pi)=\Sp_2(\chi)=\bigl(\chi\oplus\chi(1),\,N\bigr),
\]
即 $N$ 把 $\chi$ 送到 $\chi(1)$，如 1.1.4 的注 (3) 所述。
\end{enumerate}

\subsubsection{把 LLC 归结到超尖点情形的配方}

为证明局部 Langlands 对应（定理 2.13），只需证明存在唯一的一族双射
\[
\rec_n\colon\Asc_n\longrightarrow\Girr_n
\]
满足 2.13 中的全部条件，其中 $\Asc_n$ 是 $\GL_n(K)$ 超尖点表示的等构类（偏序）集，
$\Girr_n$ 是 $W_K$ 的 $n$ 维不可约 Weil--Deligne 表示的等构类集。

这是因为：一旦对所有 $m$ 在 $\Asc_m$ 上有了 $\rec$，对 $\GL_n$ 的不可约容许表示
$\pi=Q(\Delta_1,\dots,\Delta_r)$（Zelevinsky 段 $\Delta_i=\Delta_i(\pi_i,m_i)$，
$n=\sum n_im_i$），就可以定义
\[
\rec_n(\pi)=\bigoplus_{i=1}^{r}\Sp_{m_i}\bigl(\rec_{n_i}(\pi_i)\bigr).
\]
注意 $\rec_n(\pi)=(\rho,N)$ 的底层表示 $\rho$ 正是直和
\[
\rho=\bigoplus_{1\le i\le r}\ \bigoplus_{j=0}^{m_i-1}\rec_{n_i}(\pi_i)(j).
\]
可以直接验证这一构造满足 2.13 的全部要求 (1)--(5)。

\subsection{平方可积、缓增与一般表示}

我们引入如下表示（的包含链）：
\[
\text{超尖点}\subset\text{平方可积}\subset\text{缓增}\subset\text{一般（generic）}
\subset\text{容许}.
\]

\subsubsection{本质 $L^2$ 与本质 $L^{2+\epsilon}$}

\begin{definition}
设 $G^{0}=\{g\in G=\GL_n(K)\mid |\det(g)|_K=1\}$。
\begin{enumerate}[label=(\arabic*),leftmargin=2.2em]
\item 对任意 $r\in\mathbb{R}_{>0}$，容许表示 $(\pi,V)$ 称为\emph{本质 $L^r$}，
      若对任意 $v\in V$、$\lambda\in V^{\vee}$，矩阵系数 $c_{v,\lambda}$ 在 $G^{0}$ 上属于 $L^r$，即
\[
\int_{G^{0}}|c_{v,\lambda}|^{r}\,dg
\]
存在。
\item $(\pi,V)\in\Repadm_G$ 称为 $L^r$，若它本质 $L^r$ 且中心特征标 $\omega_{\pi}$ 酉。
\item $(\pi,V)\in\Repadm_G$ 称为\emph{本质平方可积}（相应地，\emph{平方可积}），
      若它本质 $L^2$（相应地，$L^2$）。
\end{enumerate}
\end{definition}

\begin{remark}
对 $(\pi,V)\in\Repadm_G$，若 $\pi$ 本质 $L^r$，则它对所有 $r'\ge r$ 也都本质 $L^{r'}$。
\end{remark}

\begin{remark}[Casselman 1980]
对任一中心特征标为 $\omega_{\pi}$ 的光滑表示 $(\pi,V)\in\Repinf_G$，
存在（唯一的？）正实值拟特征标 $\chi$，使 $\omega_{\chi\pi}$ 酉。
这样我们总可以通过扭转，把本质 $L^r$ 表示变成 $L^r$ 表示。
\end{remark}

\begin{theorem}[Zelevinsky]
\leavevmode
\begin{enumerate}[label=(\arabic*),leftmargin=2.2em]
\item $\GL_n(K)$ 的不可约容许表示 $\pi$ 本质平方可积，当且仅当它是（单个）Zelevinsky 段
      的 $Q(\Delta)$。
\item 上述本质 $L^2$ 表示 $Q(\Delta)$ 平方可积，当且仅当 $\rho\bigl(\frac{m-1}{2}\bigr)$
      的中心特征标酉，其中 $\Delta=\Delta(\sigma,m)$ 长度为 $m$。
\end{enumerate}
\end{theorem}

\subsubsection{本质缓增与缓增}

我们需要一个比（本质）平方可积稍弱的概念。

\begin{definition}
表示称为\emph{本质缓增}，若对所有 $\epsilon>0$ 它都是本质 $L^{2+\epsilon}$；
若进一步其中心特征标酉，则称为\emph{缓增}。
\end{definition}

\begin{theorem}[Zelevinsky]
$\GL_n(K)$ 的不可约容许表示 $\pi=Q(\Delta_1,\dots,\Delta_r)$ 缓增，
当且仅当每个 $Q(\Delta_i)$ 平方可积。
\end{theorem}

\begin{corollary}
若 $\pi=Q(\Delta_1,\dots,\Delta_r)$ 缓增，则任两个区间 $\Delta_i$ 与 $\Delta_j$ 都不相链；
特别地
\[
\pi\cong Q(\Delta_1)\times\cdots\times Q(\Delta_r).
\]
\end{corollary}

\begin{proof}
设 $\Delta_i=\Delta(\sigma_i,m_i)$。因 $\pi=Q(\Delta_1,\dots,\Delta_r)$ 缓增，
$\sigma_i\bigl(\frac{m_i-1}{2}\bigr)$ 的特征标酉。而若 $\sigma$ 的特征标酉，
则对任一非零 $s\in\mathbb{R}$，$\omega_{\sigma(s)}$ 不酉。
\end{proof}

\begin{remark}
由上述推论，缓增表示恰是那些每个 $Q(\Delta_i)$ 的中心特征标酉的一般表示。
用 Zelevinsky 段来分类本质缓增表示似乎（？）是困难的。
\end{remark}

\subsubsection{Whittaker 模型与一般表示\en{11}}

设 $\psi\colon K\to\mathbb{C}^{\times}$ 是 $K$ 上的非平凡加法拟特征标（如 1.1.10 小节），
$n(\psi)$ 是使 $\psi(\varpi^{-n}\mathcal{O}_K)=1$ 的最大整数（同前）。
设 $U_n\subset\GL_n(K)$ 是严格上三角矩阵构成的子群。定义 $U_n$ 的一维表示 $\theta_{\psi}$：
\[
\theta_{\psi}(u_{ij})=\psi(u_{12}+\cdots+u_{n-1,n}).
\]

\begin{definition}
$\GL_n(K)$ 的不可约容许表示 $\pi$ 称为\emph{一般}（generic），若
$\Hom_{U_n}(\pi|_{U_n},\theta_{\psi})\neq 0$。
\end{definition}

\begin{definition}
设 $(\pi,V)$ 是 $G=\GL_n(K)$ 的不可约容许一般表示，$\pi$ 关于 $\psi$ 的
\emph{Whittaker 模型}定义如下：取非零元素
$\lambda\in\Hom_{U_n}(\pi|_{U_n},\theta_{\psi})$，令
\[
W_{\psi}:=\{f\colon G\to\mathbb{C}\mid f(ug)=\theta_{\psi}(u)f(g)\ 
\ \forall g\in G,\ u\in U_n\}.
\]
$G$ 通过右平移作用在 $W_{\psi}$ 上。定义（单射的）$G$-同态
\[
V\longrightarrow W_{\psi},\qquad v\longmapsto\bigl(g\mapsto\lambda(\pi(g)v)\bigr).
\]
其像 $W_{\pi,\psi}$（同构于 $(\pi,V)$）称为 $\pi$ 的 Whittaker 模型。
\end{definition}

\begin{lemma}
\leavevmode
\begin{enumerate}[label=(\arabic*),leftmargin=2.2em]
\item 「一般」的定义与 $\psi$ 的选取无关。
\item $\pi$ 一般，当且仅当 $\pi^{\vee}$ 一般。
\item 若 $\chi\colon K^{\times}\to\mathbb{C}^{\times}$ 是任一拟特征标，则 $\pi$ 一般
      当且仅当 $\chi\pi$ 一般。
\end{enumerate}
\end{lemma}

\begin{theorem}[Zelevinsky]
$\GL_n(K)$ 的不可约容许表示 $\pi=Q(\Delta_1,\dots,\Delta_r)$ 一般，
当且仅当任两个 Zelevinsky 段 $\Delta_i$ 与 $\Delta_j$ 都不相链；由定理 3.1 第 (5) 条，
这又等价于
\[
\pi=Q(\Delta_1)\times\cdots\times Q(\Delta_r).
\]
特别地，推论 3.4 可以重新表述为：每个不可约（本质）缓增表示 $(\pi,V)$ 都是一般的。
\end{theorem}

\begin{remark}[Shalika]
一般表示大体上就是整体应用中我们所关心的表示。设 $F$ 是数域，则任一不可约容许
尖点表示 $\pi=\otimes'_v\pi_v$ 的所有局部分量 $\pi_v$ 都是一般的。
\end{remark}

\subsection{$K$-球面表示}

设 $K\subset G=\GL_n(K)$ 是紧开子群，不可约容许表示 $\pi$ 称为 \emph{$K$-球面}的，
若 $\pi^{K}\neq 0$；若对极大紧开子群 $K_0=\GL_n(\mathcal{O}_K)$ 是 $K_0$-球面的，
则称 $\pi$ 为\emph{球面的}。本小节考虑球面表示与 Iwahori-球面表示。

\subsubsection{$\GL_n(K)$ 的结构}

记号：
\begin{itemize}[leftmargin=1.6em]
\item 设 $B\subset G=\GL_n(K)$ 是标准 Borel 子群，$U\subset B$ 是严格上三角矩阵，
      $U^{-}$ 是严格下三角矩阵。设 $T\subset B$ 是对角环面（$B$ 相应的 Levi）。
\item $K_0=\GL_n(\mathcal{O}_K)$ 是极大紧子群。
\item 设 $T_0:=T\cap K_0\cong(\mathcal{O}_K^{\times})^{n}$。
\item 设 $B_0\subset K_0$ 是 Iwahori 子群，即标准 Borel $B(k)\subset\GL_n(k)$
      在约化映射下的原像。在本文中我们交替使用 Iw $=B_0$。
\item 设 $N=N_G(T)\subset G$ 是（广义）置换矩阵群，$N_K=N\cap K$。
      我们把 $N_K/T_0$ 与 Weyl 群 $W$ 等同。
\item 设 $W_1:=N/T_0$，有时称为\emph{广义 Weyl 群}。
\end{itemize}

回忆下述分解。

\begin{lemma}
\leavevmode
\begin{enumerate}[label=(\arabic*),leftmargin=2.2em]
\item $\displaystyle K_0=\bigsqcup_{w\in W}B_0wB_0$。
\item $G$ 的 Iwasawa 分解：$G=BK_0$。
\item Cartan 分解：设 $T^{-}\subset T$ 是对角矩阵
      $\mathrm{Diag}(\varpi_K^{r_1},\varpi_K^{r_2},\dots,\varpi_K^{r_n})$（$r_1\ge r_2\ge\cdots\ge r_n$）
      构成的子集，则
\[
G=\bigsqcup_{\lambda\in T^{-}/T^{0}}K_0\lambda K_0 .
\]
\item Bruhat--Tits 分解：仍记 $W_1=N/T_0$，则
\[
G=\bigsqcup_{w_1\in W_1}B_0w_1B_0 .
\]
\item Iwahori 分解：$B_0=(B_0\cap U^{-})\cdot T_0\cdot(B_0\cap U)$。
\end{enumerate}
\end{lemma}

讨论 Hecke 算子时，把 $K_0\lambda K_0$ 写成一族互不相交的 $\lambda'K_0$ 之并
也有用；初稿中我们跳过这一点。

\subsubsection{非分歧主序列}

底本此处为空（仅标题，作者原拟在二稿补写）。

\subsubsection{非分歧表示}

$G=\GL_n(K)$ 的每个不可约非分歧表示 $\pi$ 都形如 $Q(\chi_1,\dots,\chi_n)$，
其中 $\chi_1,\dots,\chi_n$ 是 $K^{\times}$ 的非分歧拟特征标。

\begin{example}
考虑 $\GL_2(K)$ 的情形。上述断言说明 $\GL_2(K)$ 的不可约非分歧表示是下列之一：
\begin{enumerate}[label=(\arabic*),leftmargin=2.2em]
\item $\pi=\chi_1\times\chi_2$，其中 $\chi_1/\chi_2\neq|\cdot|^{\pm 1}$，
      是 $\chi_1,\chi_2$ 为非分歧拟特征标的主序列表示；
\item $\pi=Q(\chi|\cdot|^{\frac12},\chi|\cdot|^{-\frac12})=\triv\otimes\chi$，
      其中 $\chi$ 是非分歧拟特征标。
\end{enumerate}
\end{example}

\subsubsection{Iwahori-不动向量}

底本此处写作「to be added in second draft」（留待二稿补写）。

\subsection{容许表示的导子与新向量}

底本此节只有标题，正文未写。

\subsection{深度 $0$ 的表示}

我们把这一讨论总结为下图。这里只关心自守一侧的不可约容许表示，以及
Weil--Deligne 表示 $\rec_n(\pi)=(\rho,N)$（或相应的 $l$-进表示 $r_n(\pi)$，
使 $\rec_n(\pi)=\mathrm{WD}(r_n(\pi))$）。

\begin{center}
\renewcommand{\arraystretch}{1.35}
\begin{tabular}{@{}r@{\qquad$\Longleftrightarrow$\qquad}l@{}}
$\pi$ 非分歧 & $\pi^{K}\neq 0$\\
$r_n(\pi)$ 非分歧 & $\rho|_{I_K}=0$ 且 $N=0$\\
$\pi$ 是 Iwahori-球面的 & $\pi^{\mathrm{Iw}}\neq 0$\\
$\rec_n(\pi)$ 为 $I_K$-球面的 & $\rho|_{I_K}=0$\\
$\pi$ 的深度为 $0$ & $r_n(\pi)$ 驯分歧
\end{tabular}
\end{center}

\subsubsection{导子与深度}

底本此处只有标题，正文未写。

\subsection{$L$-因子与 $\epsilon$-因子}

底本此处写作「To be completed in second draft」（留待二稿补完）。

\subsubsection{一般表示的 $\epsilon$-因子}

我们把这一讨论留到将来的草稿，但先记下如下几点。
表示论一侧的 $\epsilon$-因子，是以超尖点情形的 $\epsilon$-因子为基础、
依托 Bernstein--Zelevinsky 分类来定义的。对一对表示 $\pi,\pi'$，我们将定义
$\epsilon(\pi\times\pi',s,\psi)$。

\subsubsection{归纳定义}

底本此处只有标题，正文未写。

\subsubsection{$\epsilon$-因子与导子}

设 $\pi$ 是 $\GL_n(K)$ 的光滑不可约表示。当 $\pi'\colon K^{\times}\to\mathbb{C}^{\times}$
取平凡特征标时，我们用 $L(\pi,s)$ 与 $\epsilon(\pi,s,\psi)$ 记相应的 $L$-因子与
$\epsilon$-因子。仍记 $q=\#k$。由 JPPS 的定理（及其勘误），有
\[
\epsilon(\pi,s,\psi)=q^{-\left(c(\pi)+n\cdot n(\psi)\right)\cdot s}\cdot
\epsilon(\pi,0,\psi),
\]
其中 $c(\pi)$ 是 3.4 小节定义的 $\pi$ 的导子。（特别地，文献中 $\pi$ 的导子的
两个定义——其一是 3.4 小节的定义，其二是由 $L(\pi,s)$ 的函数方程给出的定义——
彼此一致。）

\subsubsection{导子被 LLC 保持}

现在有下述平凡的观察。

\begin{lemma}
在局部 Langlands 对应之下，有
\[
c(\pi)=c\bigl(\rec_n(\pi)\bigr).
\]
换言之，导子被保持。
\end{lemma}

\begin{proof}
取 $\pi'$ 为平凡特征标，由 $\epsilon$-因子的相容性得
$\epsilon(\pi,\psi,s)=\epsilon(\rec_n(\pi),\psi,s)$，特别地，由推论 1.4 及上述讨论，
\[
q^{-\left(c(\pi)+n\cdot n(\psi)\right)\cdot s}
=q^{-\left(c(\rec_n(\pi))+n\cdot n(\psi)\right)\cdot s}.
\]
\end{proof}

\begin{keytakeaways}
\begin{itemize}[leftmargin=1.5em]
\item Bernstein--Zelevinsky 分类：$\GL_n$ 的不可约容许表示唯一地写成
      $Q(\Delta_1,\dots,\Delta_r)$，段族在 $(*)$-置换下唯一；$\pi(\Delta)$ 的合成列长度 $2^{m-1}$。
\item 该分类给出 $R_K=\bigoplus_n R_n(K)$ 的分次交换环结构与对合 $Q\mapsto Z$，
      并给出基本构件：超尖点（单段且段长 $1$ 时）、本质平方可积（单段）、
      缓增（各段平方可积）、一般（各段互不相链）。
\item LLC 的证明可以归结为超尖点情形：$\rec_n(\pi)=\bigoplus_i\Sp_{m_i}(\rec_{n_i}(\pi_i))$，
      其余性质自动成立；这就是 Henniart 论证的骨架（3.1.7 与 2.3.1 合用）。
\item $\GL_2$ 的四类表示（超尖点、主序列、一维、Steinberg 扭转）分别对应
      $(\rho,0)$ 不可约、$(\chi_1\oplus\chi_2,0)$、$(\chi\oplus\chi(1),0)$、$\Sp_2(\chi)$。
\item 导子被 LLC 保持：$c(\pi)=c(\rec_n(\pi))$，由 $\epsilon$-因子的相容性与
      JPPS 公式 $q^{-(c(\pi)+n\,n(\psi))s}$ 直接得到。
\end{itemize}
\end{keytakeaways}
